\section{Revisión de literatura}
\label{sec:literatura}

\subsection{Perfil y planeza: lo que la laminación en caliente le entrega al Skinpass}
En la laminación en caliente, el espesor y la corona son los indicadores de calidad que más pesan en la
rentabilidad~\cite{peng2014}, y la literatura de procesos los trata como resultados que el molino debe controlar. La
planeza, en cambio, depende de la distribución de alargamientos a lo ancho: una banda cuya corona no guarda
proporción con la del claro entre rodillos se alarga de manera desigual y acumula esfuerzos residuales que más tarde
aparecen como ondas o pandeo, y la planeza no uniforme con que llega la banda es una dificultad reconocida de la
laminación de temple~\cite{lenard2014}. Esa es la conexión física que motiva este trabajo: la corona con que
sale el rollo del MC4 debería importar para el Skinpass, porque fija cuánta elongación diferencial tendría que
compensar la pasada de temple. Pero en una planta la corona no es una variable libre. Depende del pedido (espesor y
ancho objetivo, grado), del estado de los rodillos del molino y del periodo de producción, de modo que cualquier
comparación entre rollos con distinta corona compara también pedidos y momentos distintos. La literatura de
laminación explica el mecanismo; no dice cuánto de ese mecanismo sobrevive a la mezcla de pedidos de un año real.

\subsection{La laminación de temple}
Los modelos de la laminación de temple coinciden en que es un proceso difícil de describir con la teoría clásica de
laminación plana. Con reducciones tan pequeñas, las regiones elástica y plástica de la banda son comparables, el
aplanamiento de los rodillos de trabajo domina las fuerzas y la reducción es del mismo orden que la rugosidad
superficial; por eso los modelos tradicionales predicen mal las fuerzas de separación y deben incorporar la
deformación del rodillo y fricción alta~\cite{lenard2014}. Kijima y Bay~\cite{kijima2008} ilustran la escala del
problema: con simulaciones de elementos finitos de una indentación plana, estudian cómo se transfiere la rugosidad del
rodillo a la lámina y cuánto se alarga ésta bajo carga normal pura, en un régimen de fricción seca y deformación muy
inhomogénea. Otros modelos tratan el molino de temple como un sistema de control de múltiples entradas y salidas, con
las tensiones, la fuerza y la velocidad como variables acopladas~\cite{lenard2014}. En conjunto, esta literatura
explica qué pasa dentro de una pasada y cómo controlarla; su unidad es el contacto rodillo--lámina o la línea en
estado estable, no la población de rollos de una planta. No responde, ni pretende hacerlo, cuánto difiere entre
productos la corrección que necesita el Skinpass cuando opera con una sola receta.

\subsection{Datos de proceso y datos de ocasión}
Los registros de planta son datos que nadie diseñó para contestar preguntas causales. Box~\cite{box1966} lo advirtió
para la regresión con datos no planeados: el término de error contiene variables latentes que cambian junto con las
variables observadas, y el coeficiente estimado mezcla el efecto buscado con el de esas variables; para saber qué
pasa cuando se interviene un sistema hay que intervenirlo. La literatura de monitoreo basado en datos en laminación
en caliente aprovecha bien esos registros para vigilar indicadores~\cite{peng2014}, porque su objetivo es detectar
desviaciones, no estimar qué pasaría con otra configuración. Hernán, Hsu y Healy~\cite{hernan2019} ordenan esta
distinción en tres tareas: describir, predecir y predecir contrafactuales; solo la tercera contesta "¿qué pasaría si
cambiamos la receta?", y exige supuestos que los datos de ocasión rara vez cumplen. Nuestra pregunta se ubica a
propósito en las dos primeras.

\subsection{Inferencia con datos observacionales}
Cuando se quiere ir más allá de la descripción, la literatura de inferencia observacional pide tres cosas. Primero,
comparar unidades comparables: ajustar por las variables que determinan la exposición y verificar que haya
traslape entre grupos antes de estimar~\cite{rosenbaum1983}. Segundo, que la incertidumbre respete la estructura de
los datos: con observaciones agrupadas, los errores estándar ingenuos subestiman la incertidumbre y deben agruparse
al nivel en que se asigna la variación~\cite{cameron2015}, y las fuentes jerárquicas de variación se describen mejor
con modelos multinivel que con un solo término de error~\cite{gelman2007}. Tercero, validar sin verdad de terreno:
controles negativos que no deberían responder a la exposición~\cite{lipsitch2010}, análisis de multiverso que muestren
juntas las alternativas razonables de depuración~\cite{steegen2016} y análisis de sensibilidad que digan qué tan
fuerte tendría que ser un confusor no medido para anular la conclusión, idealmente contra un confusor observado de
referencia~\cite{cinelli2020,vanderweele2017}. Cuando se examinan muchos grupos, el control de la tasa de
descubrimientos falsos~\cite{benjamini1995} y la agregación parcial~\cite{gelman2012} limitan los hallazgos espurios,
y elegir el grupo con el mejor resultado sobrestima su efecto: la maldición del optimizador~\cite{smith2006}. Estas
herramientas se desarrollaron sobre todo en epidemiología y economía; su aplicación a la cadena de laminación es
escasa, y es lo que este trabajo intenta.

\subsection{Aprendizaje automático en manufactura}
El aprendizaje automático promete extraer patrones de procesos complejos y ya tiene aplicaciones exitosas en
manufactura, pero su adopción enfrenta dos retos comunes: conseguir datos relevantes, cuya disponibilidad, calidad y
composición condicionan el desempeño, e interpretar los resultados~\cite{wuest2016}. Dos riesgos son específicos de nuestro caso. El primero es la fuga de
información: variables posteriores al resultado o etiquetas generadas por reglas sobre las mismas mediciones inflan
el desempeño y producen modelos inútiles en operación~\cite{kaufman2012,kapoor2023}. El segundo es la validación: si
rollos del mismo programa quedan a la vez en entrenamiento y prueba, el modelo aprende el programa y no el rollo, y
el desempeño reportado es optimista; las particiones deben respetar el tiempo y los grupos~\cite{roberts2017}. Aun
evitando ambos riesgos, un buen modelo predictivo no es un modelo explicativo: la importancia de una variable dice
qué usa el modelo para predecir, no qué pasaría al intervenirla~\cite{shmueli2010}.

\subsection{Síntesis y vacío}
\label{sec:vacio}
Las cinco líneas coinciden en algo que rara vez se dice junto. La física de la laminación explica por qué la forma
de entrada debería importar al Skinpass~\cite{lenard2014}, pero la literatura que revisamos la estudia dentro de una
pasada y no con la población de rollos de una planta~\cite{kijima2008}. Los datos de planta permiten describir esa
población, pero no contestar por sí solos qué receta sería mejor~\cite{box1966,hernan2019}.

Eso responde por qué la optimización de parámetros con aprendizaje automático no resuelve por sí sola el problema de
Ternium. Un modelo que predice la calidad a partir del producto y de la configuración, y que después se optimiza,
aprende qué configuraciones tuvieron los rollos que salieron bien bajo la regla histórica de asignación, no qué pasaría
si se cambiara la receta~\cite{hernan2019,shmueli2010}. Con una receta por tipo de material, casi no hay variación de
configuración dentro de cada producto, de modo que el modelo extrapola fuera del soporte observado. El bending real
lo fija un lazo que reacciona a la forma del rollo, lo que lo convierte en mediador y no en palanca. Y sin planeza de
salida no hay un resultado de forma que optimizar.

El vacío concreto que atendemos es previo a cualquier optimización: no encontramos, en la literatura revisada, una
medición de cuánto de la corrección que aplica el lazo de bending del Skinpass se explica por el producto y cuánto por
el programa o el periodo, a partir de datos ligados entre el molino caliente y el de temple. Esa cantidad acota lo que
una receta por producto podría absorber, y es la que este artículo estima.
